\section{Introducción}


La manipulación de imágenes mediante técnicas y herramientas de inteligencia artificial ha aumentado y se ha vuelto más sofisticada y accesible para el público general. Si bien esto representa un avance en el campo de la visión por computadora, también genera una brecha técnica que dificulta distinguir lo alterado de lo original. Un caso particular son las manipulaciones faciales, conocidas como \textit{deepfakes}, las cuales pueden alcanzar un alto nivel de realismo porque conservan coherencia visual al modificar regiones específicas del rostro. Esta situación exige métodos que no solo determinen una clasificación general, sino que también localicen espacialmente la región alterada.
\subsection{Contexto del problema}

En el ámbito social, el aumento de los \textit{deepfakes} y la capacidad de sustituir rostros en imágenes o videos representan riesgos como difamación personal, creación de evidencia visual falsa, suplantación de identidad y desinformación a escala masiva.

En áreas como la política y los medios de comunicación, esto genera riesgos para la confiabilidad e integridad de la información, dado que la veracidad del contenido es difícil de avalar únicamente con medios de detección global. El alcance de estas prácticas puede ir de lo individual a lo colectivo, pues compromete contextos donde la evidencia audiovisual resulta decisiva.

\subsection{Planteamiento del problema}

Desde la perspectiva técnica, el problema se centra en que, mientras mayor realismo alcanza una manipulación, menos evidentes pueden ser las características o patrones visuales que la delatan. Muchos sistemas de detección se concentran en una clasificación binaria global, donde asignan una clase real o falsa a la imagen completa. Este enfoque puede ser útil cuando la manipulación afecta gran parte de la imagen; sin embargo, una alteración facial suele estar delimitada a regiones específicas, lo que puede reducir la sensibilidad del detector y producir falsos negativos.


\subsection{Justificación}

% Explicar por qué tiene sentido usar segmentación semántica y redes convolucionales.
% Justificar U-Net por su uso en localización píxel a píxel.
Por este motivo, el proyecto propone transitar de una clasificación global hacia la segmentación de regiones, exigiendo que el sistema delimite e identifique las áreas del rostro que han sido intervenidas. Para lograr esto, las redes neuronales convolucionales enfocadas en segmentación, como U-Net, resultan adecuadas, ya que su capacidad de evaluar píxel por píxel permite detectar anomalías que los modelos globales pueden omitir por la proporción de la región manipulada respecto a la imagen completa. Además, la segmentación binaria permite generar evidencia visual de valor interpretativo, ya que no se limita a decidir si una imagen es falsa, sino que resalta la región intervenida para guiar su inspección.

\subsection{Objetivo general}

Desarrollar una red neuronal para la detección y localización de manipulaciones faciales (deepfakes) mediante un modelo de segmentación semántica basado en redes neuronales convolucionales (arquitectura U-Net), orientado al análisis forense de imágenes.

\subsection{Objetivos específicos}

\begin{itemize}
    \item Generar un conjunto de datos (dataset) conformado por 50,000 muestras, estandarizando las imágenes para incluir variaciones de contexto y etiquetado mediante máscaras binarias que delimiten la región facial manipulada.
    \item Implementar una arquitectura basada en U-Net para la detección de inconsistencias en texturas, iluminación y bordes, capaz de generar una salida de segmentación.
    \item Entrenar el modelo utilizando el dataset generado, enfocándose en la generalización frente a técnicas de reemplazo y alteraciones localizadas en el rostro.
    \item Evaluar el desempeño del sistema con métricas de segmentación semántica y generar visualizaciones (\textit{overlays}) que resalten las regiones modificadas para facilitar la interpretación forense de los resultados.
\end{itemize}

\subsection{Alcance del proyecto}
El presente proyecto se diseñó para el análisis offline de una imagen exclusivamente dentro del contexto semántico de un rostro; se omiten cuello, orejas y cabello, y no se requiere una referencia del sujeto en un estado no manipulado. El sistema parte de características de textura, discontinuidades y bordes que el modelo pueda extraer para generar una representación aproximada de la alteración facial. Todo esto se realiza a partir del entrenamiento de la red con un conjunto de datos especializado para la tarea.
% Delimitar qué sí hace y qué no hace el proyecto.